\section{From holomorphic accuracy to logarithmic accuracy}\label{sec:logarithms}

Uniform approximation of a holomorphic function does not give a
pointwise logarithm comparison near its zeros. The appropriate local
statement uses an integral norm, a nonvanishing anchor, and a larger
disk. Multiplicity is then paid through a zero-count exponent.
The following constants are deliberately explicit rather than
optimized.

Fix once and for all
\begin{equation}\label{eq:radii}
 R=\frac92,\qquad \rho_*=\frac{13}3,\qquad r=\frac{25}6,
 \qquad B=R+r=\frac{26}3,\qquad
 \frac{\rho_*+r}{\rho_*-r}=51 .
\end{equation}
Write $D_t=D(0,t)$ in this section.

\begin{theorem}[Local logarithm comparison]\label{thm:log}
Suppose $f,g$ are holomorphic on a neighborhood of $\overline D_R$ and
\[
 |f|,|g|\le1,\qquad |f(0)|,|g(0)|\ge m,\qquad
 0<m\le1,\qquad \norm{f-g}_{L^\infty(D_R)}\le\varepsilon\le1 .
\]
Set
\begin{equation}\label{eq:local-constants}
 \begin{gathered}
 A=\log(1/m),\qquad
 n=\max\!\left(1,\left\lceil\frac{A}{\log(27/26)}\right\rceil\right),
 \quad \alpha=\frac2{n+2},\\
 C(A,n)=\pi r^2+\pi B^2n^2\ee^{102A/n}2^{2/n}.
 \end{gathered}
\end{equation}
Then, with locally integrable logarithms at zeros,
\begin{equation}\label{eq:log-transfer}
 \int_{D_r}\left|\log|f|-\log|g|\right|\dd A
 \le C(A,n)\varepsilon^\alpha .
\end{equation}
Each function has at most $n$ zeros in $D_{\rho_*}$, counted with
multiplicity. The same statements hold after translating the disks.
\end{theorem}

\begin{proof}
We give the zero count, sublevel estimate and comparison separately.
The classical inputs are Jensen's formula, the maximum principle and
the disk Harnack inequality; see
\cite{Tao2020LocalZeros,Tao2016Harnack}.

\emph{Zero count and factor removal.}
If $a_1,\ldots,a_k$ are the zeros of $f$ in $D_{\rho_*}$, repeated
according to multiplicity, Jensen's formula and $|f|\le1$ give
\[
 k\log(R/\rho_*)\le
 \sum_{|a|<R}\log(R/|a|)
 \le-\log|f(0)|\le A .
\]
The formula follows at a boundary zero by an interior-radius limit.
Thus $k\le n$, and no $a_\ell$ is zero. Remove these factors by
\[
 b_\ell(z)=\frac{R(z-a_\ell)}{R^2-\overline{a_\ell}z},\qquad
 h(z)=\frac{f(z)}{\prod_{\ell=1}^k b_\ell(z)}.
\]
The apparent poles are removable. Each $b_\ell$ has modulus one on
$|z|=R$, so $|h|\le1$ by the maximum principle.
Moreover $h$ is zero-free in $D_{\rho_*}$ and
$|h(0)|\ge|f(0)|\ge m$. Harnack applied to the nonnegative harmonic
function $-\log|h|$ yields
$-\log|h(z)|\le51A$ on $D_r$. Zeros on the circle of radius
$\rho_*$ cause no problem: apply Harnack on smaller circles and
let their radii increase to $\rho_*$.
For $|z|\le r$,
\begin{equation}\label{eq:blaschke-lower}
 |b_\ell(z)|\ge\frac{|z-a_\ell|}{B},\qquad
 |f(z)|\ge\ee^{-51A}\prod_{\ell=1}^k\frac{|z-a_\ell|}{B}.
\end{equation}
The first inequality follows from
$|R^2-\overline{a_\ell}z|\le R(R+r)$; each distance ratio is
at most one.

\emph{Sublevel areas.}
If $q=\ee^{51A}t<1$ and $k>0$, a point with $|f(z)|<t$ must lie
in one of the $k$ disks of radius $Bq^{1/k}$ around the $a_\ell$.
Their union has area at most
$k\pi B^2q^{2/k}\le n\pi B^2q^{2/n}$.
For $k=0$ the sublevel set is empty. If $q\ge1$, the same upper
bound exceeds the area of $D_r$. Hence for every $t>0$,
\begin{equation}\label{eq:sublevel}
 \bigl|\{z\in D_r:|f(z)|<t\}\bigr|
 \le L(A,n)t^{2/n},\qquad
 L(A,n)=n\pi B^2\ee^{102A/n}.
\end{equation}
Layer cake, applied to $\log(t/|f|)_+$, now gives
\begin{equation}\label{eq:layer-cake}
 \int_{\{|f|<t\}\cap D_r}\log\frac{t}{|f|}\dd A
 \le\int_0^\infty L(A,n)(t\ee^{-v})^{2/n}\dd v
 =\frac{\pi B^2n^2}{2}\ee^{102A/n}t^{2/n}.
\end{equation}
The same estimates hold for $g$.

\emph{Comparison.}
For $\varepsilon>0$ choose
$\tau=\varepsilon^{n/(n+2)}$, so $\varepsilon\le\tau$.
Where both moduli are at least $\tau$, the real logarithm is
$1/\tau$-Lipschitz, and the integral of the log difference is at
most $\pi r^2\varepsilon/\tau$.
On the complement the smaller modulus is less than $\tau$ and the
larger is less than $2\tau$. There
\[
 |\log|f|-\log|g||
 \le\log(2\tau/|f|)+\log(2\tau/|g|).
\]
Extend each nonnegative term to its full sublevel set and apply
\eqref{eq:layer-cake} twice. The resulting bound is
\[
 \pi r^2\frac{\varepsilon}{\tau}
 +\pi B^2n^2\ee^{102A/n}(2\tau)^{2/n}
 =C(A,n)\varepsilon^{2/(n+2)}.
\]
When $\varepsilon=0$ the assertion is immediate.
\end{proof}

\begin{corollary}[Compact zero-statistic transfer]\label{cor:weak-transfer}
Under the hypotheses of Theorem~\ref{thm:log}, for any
$\varphi\in C_c^{1,1}(D_r)$,
\begin{equation}\label{eq:weak-transfer}
 |\nu_f(\varphi)-\nu_g(\varphi)|
 \le \frac{\norm{\Delta\varphi}_\infty}{2\pi}
              C(A,n)\varepsilon^{2/(n+2)} .
\end{equation}
No simplicity, matching of roots, or lower bound for the modulus
throughout the observation region is required.
\end{corollary}
\begin{proof}
Subtract \eqref{eq:weak-onepoint} for $f$ and $g$ and apply
\eqref{eq:log-transfer}.
\end{proof}

We will also use an elementary bound on the individual logarithms:
\begin{equation}\label{eq:log-mass}
 \norm{\log|f|}_{L^1(D_r)},\
 \norm{\log|g|}_{L^1(D_r)}
 \le \pi r^2A .
\end{equation}
Both logarithms are nonpositive, while Jensen's formula says that each
circular mean is at least its value at the center, hence at least
$-A$. Integrating those means in the radius proves
\eqref{eq:log-mass}. This estimate concerns the size of the logarithms;
\eqref{eq:log-transfer} concerns their difference and is the useful
small quantity.
